% !TEX program = lualatex
% Transparencias de Fundamentos de las Matemáticas I
% Clase 01: Lenguaje matemático I

\documentclass[aspectratio=169,11pt]{beamer}

\usepackage{fontspec}
\usepackage[spanish,es-nodecimaldot]{babel}
\usepackage{amsmath,amssymb,mathtools}
\usepackage{booktabs}
\usepackage{csquotes}
\usepackage{xcolor}

\usetheme{Madrid}
\usecolortheme{default}
\setbeamertemplate{navigation symbols}{}
\setbeamertemplate{footline}[frame number]

\definecolor{FdMBlue}{RGB}{20,55,100}
\definecolor{FdMRed}{RGB}{130,30,30}
\definecolor{FdMGreen}{RGB}{25,100,60}

\setbeamercolor{structure}{fg=FdMBlue}
\setbeamercolor{title}{fg=white,bg=FdMBlue}
\setbeamercolor{frametitle}{fg=white,bg=FdMBlue}
\setbeamercolor{block title}{fg=white,bg=FdMBlue}
\setbeamercolor{block body}{bg=blue!3}
\setbeamercolor{alerted text}{fg=FdMRed}

\setmainfont{Latin Modern Roman}
\setsansfont{Latin Modern Sans}
\setmonofont{Latin Modern Mono}

\newcommand{\N}{\mathbb{N}}
\newcommand{\Z}{\mathbb{Z}}
\newcommand{\Q}{\mathbb{Q}}
\newcommand{\R}{\mathbb{R}}
\newcommand{\C}{\mathbb{C}}
\newcommand{\Pow}{\mathcal{P}}
\newcommand{\id}{\operatorname{id}}
\newcommand{\mcd}{\operatorname{mcd}}
\newcommand{\mcm}{\operatorname{mcm}}

\title[FdM I -- Clase 01]{Lenguaje matemático I}
\subtitle{Objetos, símbolos y significado}
\author{Antonio Falcó Montesinos}
\institute{Universidad CEU Cardenal Herrera}
\date{Fundamentos de las Matemáticas I}

\begin{document}

\begin{frame}
  \titlepage
\end{frame}

\begin{frame}{Objetivos de la clase}
\begin{itemize}
  \item Distinguir entre lenguaje natural y lenguaje matemático.
  \item Reconocer objetos, símbolos, variables y constantes.
  \item Aplicar el principio de definición previa.
\end{itemize}
\end{frame}

\begin{frame}{Mapa de la clase}
\tableofcontents
\end{frame}

\section{Punto de partida}

\begin{frame}{Punto de partida}
\begin{block}{Por qué empezar por el lenguaje}
\begin{itemize}
  \item En matemáticas no basta con tener una intuición correcta.
  \item Una afirmación debe tener significado preciso en un contexto fijado.
  \item El curso se apoya en una regla de calidad: todo símbolo relevante debe definirse antes de utilizarse.
\end{itemize}
\end{block}
\end{frame}

\begin{frame}{Punto de partida}
\begin{block}{Lenguaje natural y lenguaje matemático}
\begin{itemize}
  \item El lenguaje natural es flexible, contextual y a veces ambiguo.
  \item El lenguaje matemático reduce ambigüedad mediante definiciones, símbolos y reglas.
  \item Un buen texto matemático combina frases completas con notación bien declarada.
\end{itemize}
\end{block}
\end{frame}

\begin{frame}{Punto de partida}
\begin{block}{Ejemplo de ambigüedad}
\begin{itemize}
  \item La frase ``todo número tiene un siguiente'' no es precisa por sí sola.
  \item Hay que indicar el conjunto de trabajo y qué significa ``siguiente''.
  \item En \(\mathbb{N}\), puede formularse como: para todo \(n\in\mathbb{N}\), existe \(m\in\mathbb{N}\) tal que \(m=n+1\).
\end{itemize}
\end{block}
\end{frame}

\section{Objetos y símbolos}

\begin{frame}{Objetos y símbolos}
\begin{block}{Objeto matemático}
\begin{itemize}
  \item Un objeto matemático es una entidad definida dentro de un contexto.
  \item Puede ser un número, conjunto, función, relación, operación o estructura.
  \item El mismo símbolo puede representar objetos distintos en problemas distintos.
\end{itemize}
\end{block}
\end{frame}

\begin{frame}{Objetos y símbolos}
\begin{block}{Símbolos con significado}
\begin{itemize}
  \item En \(x\in A\), el símbolo \(x\) representa un objeto, \(A\) un conjunto y \(\in\) la relación de pertenencia.
  \item La expresión solo está justificada si esos tres elementos han sido introducidos.
  \item La notación no sustituye a la explicación.
\end{itemize}
\end{block}
\end{frame}

\begin{frame}{Objetos y símbolos}
\begin{block}{Variables y constantes}
\begin{itemize}
  \item Una variable puede representar distintos objetos dentro de un contexto.
  \item Una constante representa un objeto fijo.
  \item Ejemplo: en ``sea \(n\in\mathbb{N}\)'', \(n\) es variable; \(0\) es una constante una vez definida.
\end{itemize}
\end{block}
\end{frame}

\section{Trabajo de aula}

\begin{frame}{Trabajo de aula}
\begin{block}{Actividad}
\begin{itemize}
  \item Localizar símbolos no definidos en un texto breve.
  \item Reescribir una frase ambigua usando contexto explícito.
  \item Separar objetos, propiedades y relaciones en una fórmula dada.
\end{itemize}
\end{block}
\end{frame}

\section{Profundización}

\begin{frame}{Profundización}
\begin{block}{Criterio de lectura}
\begin{itemize}
  \item Al leer una fórmula, conviene identificar primero los objetos que aparecen.
  \item Después se localizan las relaciones u operaciones que conectan esos objetos.
  \item Finalmente se reconstruye la frase completa en lenguaje natural.
\end{itemize}
\end{block}
\end{frame}

\begin{frame}{Profundización}
\begin{block}{Ejemplo guiado}
\begin{itemize}
  \item En \(x\in A\Rightarrow x\in B\), los símbolos \(x,A,B\) deben estar declarados.
  \item La flecha \(\Rightarrow\) expresa una implicación, no una continuación de cálculo.
  \item Una lectura posible es: si \(x\) pertenece a \(A\), entonces \(x\) pertenece a \(B\).
\end{itemize}
\end{block}
\end{frame}

\begin{frame}{Profundización}
\begin{block}{Preguntas de control}
\begin{itemize}
  \item ¿Qué símbolos aparecen en la expresión?
  \item ¿Qué tipo de objeto representa cada símbolo?
  \item ¿Hay algún símbolo que el texto use sin haber definido?
\end{itemize}
\end{block}
\end{frame}

\begin{frame}{Cierre}
\begin{itemize}
  \item Un símbolo sin definición no aporta rigor.
  \item La precisión matemática empieza antes de calcular.
  \item En la próxima clase: definiciones, enunciados, ejemplos y demostraciones.
\end{itemize}
\end{frame}

\begin{frame}{Trabajo recomendado}
\begin{itemize}
  \item Releer las secciones correspondientes del manual.
  \item Identificar definiciones, símbolos nuevos y enunciados principales.
  \item Preparar dudas y ejemplos para el seminario asociado.
\end{itemize}
\end{frame}

\end{document}
